上一单元把形状和布局安顿好了，接下来的问题是怎样把导数算出来。这一单元给的答案是一套可以机械执行的动作：给每个中间量加一个扰动，只保留一阶项，用代数规则把扰动从输入传到输出，最后把结果整理成内积的样子，梯度自己就从括号里显出来了。

以矩阵最小二乘为例。记 \(R=AX-B\)，微分是 \(\mathrm df=\trace(R^{\trans}A\,\mathrm dX)\)，整理成 \(\trace\bigl((A^{\trans}R)^{\trans}\mathrm dX\bigr)\)，读出 \(\nabla_Xf=A^{\trans}(AX-B)\)。整个过程没有出现一个下标，也没有用到任何需要背的公式表——用到的只有微分规则、迹的循环性质和一次形状核对。

三篇按阶数推进。第一篇把``写微分、整理成内积、读梯度''这三步立成程序，并用它推一个新层的前向权重梯度，顺带说明和与平均、约束变量这两处容易埋雷的地方。第二篇处理复合：链式法则给出的是一串 Jacobian 连乘，而括号打在哪一端决定了代价是 \(O(Ld^3)\) 还是 \(O(Ld^2)\)，前向模式与反向模式的分野就在这里，反向传播为什么只比前向贵一倍也在这里得到完整回答。第三篇进入二阶：Hessian 是一个双线性形式，对称性来自混合偏导可交换而那需要条件；实践中很少显式构造它，用的是 Hessian-向量积。

\subsection{本章篇目}

以下三篇按概念依赖排列；若从具体困惑进入，也可以先打开最接近的一篇，再沿篇内链接回补。

\begin{itemize}
\tightlist
\item \hyperref[sec:10-Matrix-Calculus/02-Differentials/01]{``用微分与 Frobenius 内积读取梯度''}；
\item \hyperref[sec:10-Matrix-Calculus/02-Differentials/02]{``链式法则就是线性映射复合''}；
\item \hyperref[sec:10-Matrix-Calculus/02-Differentials/03]{``Hessian 是二阶线性算子''}。
\end{itemize}

读的时候反复问两件事就够了：这一项对扰动是不是线性的，最后内积里的系数与被求导的变量是不是同形。这两个检查比记住任何一条公式都管用，也会一直用到后面的迹方法与反向传播。
